\section{Introducción}

Una unidad aritmético-lógica (ALU) reúne operaciones aritméticas,
lógicas y de comparación en un circuito cuya salida se selecciona
mediante un código de operación. Su descripción en VHDL permite
estudiar tanto el comportamiento combinacional como la organización
interna del hardware y su verificación mediante simulación.

El Laboratorio 2 de Electrónica Digital 2 propone una ALU de
dieciséis operaciones, con dos operandos sin signo de cuatro bits,
un resultado del mismo ancho y las banderas Carry y Zero. El trabajo
se divide en una descripción monolítica, Parte A, y una
reestructuración jerárquica, Parte B, que conserva la interfaz y el
comportamiento de la primera versión~\cite{guialab2}.

Esta organización permite relacionar la especificación funcional
con dos formas de describir el mismo sistema. Además de comprobar
las operaciones, interesa analizar cómo la división en módulos
facilita reutilizar un sumador y verificar cada bloque antes de
integrarlo en la ALU completa.

\subsection{Objetivos}
\begin{itemize}
\item Describir una ALU combinacional de dieciséis operaciones
utilizando el paquete \texttt{numeric\_std}.
\item Implementar las banderas de estado y verificar los casos de
borde de la aritmética de cuatro bits.
\item Reorganizar el diseño mediante componentes, reutilizando un
único sumador dentro de la unidad aritmética.
\item Verificar cada módulo y el sistema completo mediante bancos
de prueba autoverificables.
\item Contrastar las dos versiones y validar su funcionamiento en
la placa Spartan Edge Accelerator con el shield UTEC.
\end{itemize}
